% Created 2026-08-17 Mon 17:48
% Intended LaTeX compiler: pdflatex
\documentclass[11pt]{article}
\usepackage[utf8]{inputenc}
\usepackage[T1]{fontenc}
\usepackage{graphicx}
\usepackage{longtable}
\usepackage{wrapfig}
\usepackage{rotating}
\usepackage[normalem]{ulem}
\usepackage{amsmath}
\usepackage{amssymb}
\usepackage{capt-of}
\usepackage{hyperref}
\author{user}
\date{\today}
\title{}
\hypersetup{
 pdfauthor={user},
 pdftitle={},
 pdfkeywords={},
 pdfsubject={},
 pdfcreator={Emacs 30.2 (Org mode 9.7.11)}, 
 pdflang={English}}
\begin{document}

\tableofcontents

\section{Task 1}
\label{sec:org4ec4dce}
The task is to not use reverse() or split (using [::-1]) and
print an arr backwards. For this we use the "for" function and
loop the original arr backwards and add the value to the new arr.
\begin{verbatim}
num_list = [3, 6, 9, 12]
def home_made_reverse(nlist):
	newlist = []
	for i in range(len(nlist) -1, -1, -1):
		newlist.append(nlist[i])
	return newlist

modded_num_list = home_made_reverse(num_list)
print(modded_num_list)
\end{verbatim}

\phantomsection
\label{orgc234817}
\begin{verbatim}
[12, 9, 6, 3]
\end{verbatim}
\section{Task 2}
\label{sec:org99a7742}
The task is to implement a sorting algorithm, which i implemented
bubble sort.
\begin{verbatim}
def bubble_sort(arr, reverse=False):
    n = len(arr)
    # Traverse through all array elements
    for i in range(n):
        # Last i elements are already sorted, no need to check them
        for j in range(0, n-i-1):
            if (arr[j] > arr[j+1]) != reverse:
                arr[j], arr[j+1] = arr[j+1], arr[j]

arr1 = [64, 34, 25, 12, 22, 11, 90]
bubble_sort(arr1)
print("Sorted array (ascending):", arr1)

arr2 = [64, 34, 25, 12, 22, 11, 90]
bubble_sort(arr2, reverse=True)
print("Sorted array (descending):", arr2)
\end{verbatim}

\phantomsection
\label{org27ffc41}
\begin{verbatim}
Sorted array (ascending): [11, 12, 22, 25, 34, 64, 90]
Sorted array (descending): [90, 64, 34, 25, 22, 12, 11]
\end{verbatim}
\end{document}
